\section{Related Work}
\paragraph{Neural Scaling Laws for Language Models}
Neural scaling laws relate the optimal number of model parameters and amount of training data for a fixed amount of compute.  \citet{hoffmann2022training} presented the Chinchilla scaling laws for language models demonstrating that there was a linear relationship between the size of the model and the amount of training data needed. Their findings indicated that prior models such as Gopher~\citep{rae2021scaling} are severely undertrained. Recently, models such as Llama~\citep{touvron2023llama} are trained with much more data.  These scaling laws were drawn for the paradigm of single-epoch training.  Recently, \citet{muennighoff2023scaling} showed that
the marginal utility of repeated data rapidly diminishes when training for more than 4 epochs, and formulated scaling laws under repeated data.
Concurrently, \citet{xue2023repeat} showed that repeating even small fractions of the pre-training data can lead to overfitting and reduce model performance.

\paragraph{Dataset Selection}
Selecting high quality data for pre-training LLMs remains an active, high-impact, yet understudied area of research. For instance,
 GPT-2 model was pre-trained on all outbound links from Reddit, a social media platform, which received at least 3 karma~\citep{brown2020language}. This was used as a heuristic indicator that the document may be \emph{interesting, educational, or just funny.}
 Follow-up works have used other heuristics such as prioritizing documents that resemble wikipedia~\citep{gururangan-etal-2022-whose}.~\citet{rae2021scaling}
used multiple heuristic filters to remove documents, such as the absence of certain stopwords, length of the document, percentage of alphabetic characters, 
 mean word length, 
  symbol-to-word ratio, percentage of lines starting with a bullet point, or ending with an ellipsis etc. Their work highlights the intricacies of filtering out text data.
An alternative paradigm for building better datasets for training is to distill high-quality datasets.  \citet{xie2023doremi} proposed a method, DoReMi, to select the best data mixture for pre-training language models by reweighting data from various domains.
Concurrently, \citet{Abbas2023SemDeDupDL} showed that de-duplicating pre-training data can improve pre-training efficiency.
Recently several methods were proposed for automatic filtering of low-quality data for faster fine-tuning of LLMs~\citep{chen2023alpagasus,solaiman2021process,zhou2023lima}. Simultaneously, in the realm of image-language models such as CLIP~\citep{radford2021learning}, the Datacomp benchmark~\citep{gadre2023datacomp} and recent entries~\citep{maini2023t,yu2023devil} have developed approaches at filtering out low-quality subsets from pre-training datasets like LAION~\citep{schuhmann2022laion}, or from scrapes of the common crawl.


\paragraph{Data Augmentation and synthetic data}
\citet{eldan2023tinystories} showed that a synthetically generated dataset in the form of stories that toddlers can understand
allows training a small language model that can generate coherent sentences.
\citet{gunasekar2023textbooks} showed that textbook quality (synthetic) data alone helps models achieve state-of-the-art performance on reasoning and coding tasks. Similar approaches are used in concurrent work for enhancing coding and mathematical reasoning abilities while finetuning \cite{liu2023tinygsm,wei2023magicoder}.
\citet{shumailov2023model} show that training on synthetic data can actually be harmful for model performance, especially when
we do multiple rounds of pre-training an LLM and then training the next LLM on data generated by the previous one.
On the other hand, some other works have shown that such a strategy can actually be useful.
\citet{li2023self} and \citet{koksal2023longform} discuss how a model can generate instruction data and then fine-tune on its own generated data
to improve performance.
\citet{jung2023impossible} discuss how such repeated cycles of synthetic data can help train a very small paraphrase and summarization model that even outperforms GPT-3.

The vision and multimodal literatures have also seen a surge of works examining the use of synthetic data for training.
The works of \citet{bansal2023leaving,trabucco2023effective,azizi2023synthetic} have shown that using synthetic data in 
combination with real data achieves state of art model performance both in-distribution and out-of-distribution. \citet{cubuk2020randaugment} used generative models to generate image augmentations for better domain generalization. There are also multiple studies on increasing multiplicity of augmentations and their value for improving generalization \citep{choi2019faster,fort2021drawing,hoffer2020augment}.  
However, \citet{alemohammad2023self} showed that generated models trained for more than five cycles of their own generated data can undergo severe mode collapse.

